\section{Quasi-coherent sheaves and coherent sheaves}
\label{section:0.5}

\subsection{Quasi-coherent sheaves}
\label{subsection:0.5.1}

\begin{env}[5.1.1]
\label{0.5.1.1}
Let $(X,\mathcal{O}_X)$ be a ringed space and $\mathcal{F}$ an
$\mathcal{O}_X$-module. Giving a homomorphism
$u:\mathcal{O}_X\to\mathcal{F}$ of $\mathcal{O}_X$-modules is equivalent to giving
the section $s=u(1)\in\Gamma(X,\mathcal{F})$. Indeed, when $s$ is
given, for every section $t\in\Gamma(U,\mathcal{O}_X)$ one necessarily has
$u(t)=t\cdot(s|U)$; we say that $u$ is \emph{defined by the section $s$}. If
now $I$ is an arbitrary index set, consider the direct sum
sheaf $\mathcal{O}_X^{(I)}$, and for every $i\in I$, let $h_i$ be
the canonical injection of the $i$-th factor into $\mathcal{O}_X^{(I)}$; we know
that $u\mapsto(u\circ h_i)$ is an isomorphism of
$\operatorname{Hom}_{\mathcal{O}_X}(\mathcal{O}_X^{(I)},\mathcal{F})$ onto the
product
$\bigl(\operatorname{Hom}_{\mathcal{O}_X}(\mathcal{O}_X,\mathcal{F})\bigr)^I$.
There is therefore a canonical one-to-one correspondence between the homomorphisms
$u:\mathcal{O}_X^{(I)}\to\mathcal{F}$ and the \emph{families of sections
$(s_i)_{i\in I}$ of $\mathcal{F}$ over $X$}. The homomorphism $u$
corresponding to $(s_i)$ maps an element
$(a_i)\in(\Gamma(U,\mathcal{O}_X))^{(I)}$ to
$\sum_{i\in I}a_i\cdot(s_i|U)$.

We say that $\mathcal{F}$ is \emph{generated by the family $(s_i)$} if
the homomorphism $\mathcal{O}_X^{(I)}\to\mathcal{F}$ defined
\oldpage[0\textsubscript{I}]{45}
by this family is \emph{surjective} (in other words, if, for every
$x\in X$, $\mathcal{F}_x$ is an $\mathcal{O}_x$-module generated by the
$(s_i)_x$). We say that $\mathcal{F}$ is \emph{generated by its sections
over $X$} if it is generated by the family of all these sections (or
by a subfamily); in other words, if there exists a surjective homomorphism
$\mathcal{O}_X^{(I)}\to\mathcal{F}$ for a suitable $I$.

Note that an $\mathcal{O}_X$-module $\mathcal{F}$ can be such that there
exists a point $x_0\in X$ for which $\mathcal{F}|U$ is not generated by
its sections over $U$, \emph{for every open neighbourhood $U$ of
$x_0$}: it suffices to take $X=\mathbb{R}$, for $\mathcal{O}_X$ the simple
sheaf $\mathbb{Z}$, for $\mathcal{F}$ the algebraic sub-sheaf of
$\mathcal{O}_X$ such that $\mathcal{F}_0=\{0\}$, $\mathcal{F}_x=\mathbb{Z}$ for
$x\ne0$, and finally $x_0=0$: the only section of $\mathcal{F}|U$ over
$U$ is $0$ for a neighbourhood $U$ of $0$.
\end{env}

\begin{env}[5.1.2]
\label{0.5.1.2}
Let $f:X\to Y$ be a morphism of ringed spaces. If $\mathcal{F}$ is an
$\mathcal{O}_X$-module generated by its sections over $X$, then
the canonical homomorphism
$f^*(f_*(\mathcal{F}))\to\mathcal{F}$ (4.4.3.3) is \emph{surjective}:
indeed, with the notation of (5.1.1), $s_i\otimes1$ is a section of
$f^*(f_*(\mathcal{F}))$ over $X$, and its image in $\mathcal{F}$ is
$s_i$. The example of (5.1.1), where $f$ is the identity, shows that the converse of
this proposition is false in general.
\end{env}

\begin{env}[5.1.3]
\label{0.5.1.3}
An $\mathcal{O}_X$-module $\mathcal{F}$ is said to be \emph{quasi-coherent}
if, for every $x\in X$, there is an open neighbourhood $U$ of $x$ such that
$\mathcal{F}|U$ is isomorphic to the \emph{cokernel} of a homomorphism of the
form
$\mathcal{O}_X^{(I)}|U\to\mathcal{O}_X^{(J)}|U$, where $I$ and $J$ are
arbitrary index sets. It is clear that $\mathcal{O}_X$ itself is
a quasi-coherent $\mathcal{O}_X$-module, and that every direct sum of
quasi-coherent $\mathcal{O}_X$-modules is a quasi-coherent
$\mathcal{O}_X$-module. An $\mathcal{O}_X$-algebra $\mathcal{A}$ is said to be
\emph{quasi-coherent} if it is so as an $\mathcal{O}_X$-module.
\end{env}

\begin{env}[5.1.4]
\label{0.5.1.4}
Let $f:X\to Y$ be a morphism of ringed spaces. If $\mathcal{G}$ is a
quasi-coherent $\mathcal{O}_Y$-module, then $f^*(\mathcal{G})$ is a
quasi-coherent $\mathcal{O}_X$-module. Indeed, for every $x\in X$, there is
an open neighbourhood $V$ of $f(x)$ in $Y$ such that $\mathcal{G}|V$ is the
cokernel of a homomorphism
$\mathcal{O}_Y^{(I)}|V\to\mathcal{O}_Y^{(J)}|V$. If $U=f^{-1}(V)$ and
$f_U$ is the restriction of $f$ to $U$, one has
$f^*(\mathcal{G})|U=f_U^*(\mathcal{G}|V)$; since $f_U^*$ is right exact
and commutes with direct sums, $f_U^*(\mathcal{G}|V)$ is the cokernel of a
homomorphism
$\mathcal{O}_X^{(I)}|U\to\mathcal{O}_X^{(J)}|U$.
\end{env}

\subsection{Sheaves of finite type}
\label{subsection:0.5.2}

\begin{env}[5.2.1]
\label{0.5.2.1}
An $\mathcal{O}_X$-module $\mathcal{F}$ is said to be \emph{of finite type} if,
for every $x\in X$, there exists an open neighbourhood $U$ of $x$ such that
$\mathcal{F}|U$ is generated by a \emph{finite} family of sections
over $U$, or equivalently is isomorphic to a quotient sheaf of a
sheaf of the form $(\mathcal{O}_X|U)^p$ with $p$ finite. Every quotient
sheaf of a sheaf of finite type is of finite type, as is every finite direct
sum and every finite tensor product of sheaves of finite type. An
$\mathcal{O}_X$-module of finite type is not necessarily quasi-coherent,
as can be seen for the $\mathcal{O}_X$-module
$\mathcal{O}_X/\mathcal{F}$, where $\mathcal{F}$ is the example of (5.1.1). If
$\mathcal{F}$ is of finite type, then $\mathcal{F}_x$ is an
$\mathcal{O}_x$-module of finite type for every $x\in X$, but the example
(5.1.1) shows that this necessary condition is not in general
sufficient.
\end{env}

\begin{env}[5.2.2]
\label{0.5.2.2}
Let $\mathcal{F}$ be an $\mathcal{O}_X$-module \emph{of finite type}. If $s_i$
$(1\leq i\leq n)$ are sections of $\mathcal{F}$ over an open
neighbourhood $U$ of a point $x\in X$ and if the $(s_i)_x$ generate
$\mathcal{F}_x$, there exists an open neighbourhood $V\subset U$ of $x$ such that
the $(s_i)_y$ generate $\mathcal{F}_y$ for every $y\in V$ (FAC, I, 2, 12,
prop.~1). In particular, one concludes that the support of $\mathcal{F}$ is
\emph{closed}.

\oldpage[0\textsubscript{I}]{46}

Likewise, if $u:\mathcal{F}\to\mathcal{G}$ is a homomorphism such that
$u_x=0$, then there exists a neighbourhood $U$ of $x$ such that $u_y=0$ for every
$y\in U$.
\end{env}

\begin{env}[5.2.3]
\label{0.5.2.3}
Suppose $X$ \emph{quasi-compact}, and let $\mathcal{F}$, $\mathcal{G}$ be two
$\mathcal{O}_X$-modules such that $\mathcal{G}$ is \emph{of finite type},
and $u:\mathcal{F}\to\mathcal{G}$ a \emph{surjective} homomorphism. Suppose
moreover that $\mathcal{F}$ is the inductive limit of an inductive system
$(\mathcal{F}_\lambda)$ of $\mathcal{O}_X$-modules. Then there exists an index
$\mu$ such that the homomorphism $\mathcal{F}_\mu\to\mathcal{G}$ is
\emph{surjective}. Indeed, for every $x\in X$, there exists a finite system of
sections $s_i$ of $\mathcal{G}$ over an open neighbourhood $U(x)$ of $x$
such that the $(s_i)_y$ generate $\mathcal{G}_y$ for every $y\in U(x)$; there is
therefore an open neighbourhood $V(x)\subset U(x)$ of $x$ and $n$ sections $t_i$ of
$\mathcal{F}$ over $V(x)$ such that
$s_i|V(x)=u(t_i)$ for every $i$; one can moreover suppose that the $t_i$ are
the canonical images of sections of one and the same sheaf
$\mathcal{F}_{\lambda(x)}$ over $V(x)$. Now cover $X$ with
finitely many neighbourhoods $V(x_k)$, and let $\mu$ be an index greater than all
the $\lambda(x_k)$; it is clear that this index answers the question.

Still supposing $X$ quasi-compact, let $\mathcal{F}$ be an
$\mathcal{O}_X$-module of finite type generated by its sections over $X$
(5.1.1); then $\mathcal{F}$ is generated by a \emph{finite} subfamily
of these sections: indeed, it suffices to cover $X$ by finitely many
open neighbourhoods $U_k$ such that, for each $k$, there are finitely many
sections $s_{ik}$ of $\mathcal{F}$ over $X$ whose restrictions to
$U_k$ generate $\mathcal{F}|U_k$; it is clear that the $s_{ik}$ then generate
$\mathcal{F}$.
\end{env}

\begin{env}[5.2.4]
\label{0.5.2.4}
Let $f:X\to Y$ be a morphism of ringed spaces. If $\mathcal{G}$ is an
$\mathcal{O}_Y$-module of finite type, then $f^*(\mathcal{G})$ is an
$\mathcal{O}_X$-module of finite type. Indeed, for every $x\in X$, there are an
open neighbourhood $V$ of $f(x)$ in $Y$ and a surjective homomorphism
$v:\mathcal{O}_Y^p|V\to\mathcal{G}|V$. If $U=f^{-1}(V)$ and $f_U$ is the
restriction of $f$ to $U$, one has
$f^*(\mathcal{G})|U=f_U^*(\mathcal{G}|V)$; since $f_U^*$ is right exact
(4.3.1) and commutes with direct sums (4.3.2), $f_U^*(v)$ is a
surjective homomorphism
$\mathcal{O}_X^p|U\to f^*(\mathcal{G})|U$.
\end{env}

\begin{env}[5.2.5]
\label{0.5.2.5}
An $\mathcal{O}_X$-module $\mathcal{F}$ is said to \emph{admit a finite
presentation} if, for every $x\in X$, there exists an open neighbourhood $U$ of $x$ such
that $\mathcal{F}|U$ is isomorphic to a \emph{cokernel of an
$(\mathcal{O}_X|U)$-homomorphism}
$\mathcal{O}_X^p|U\to\mathcal{O}_X^q|U$, $p$ and $q$ being two integers $>0$.
Such an $\mathcal{O}_X$-module is thus of finite type and quasi-coherent. If
$f:X\to Y$ is a morphism of ringed spaces and $\mathcal{G}$ is an
$\mathcal{O}_Y$-module admitting a finite presentation, then
$f^*(\mathcal{G})$ admits a finite presentation, as the
argument of (5.1.4) shows.
\end{env}

\begin{env}[5.2.6]
\label{0.5.2.6}
Let $\mathcal{F}$ be an $\mathcal{O}_X$-module admitting a finite presentation
(5.2.5); then, for every $\mathcal{O}_X$-module $\mathcal{H}$,
the canonical functorial homomorphism
\[
  \bigl(\mathcal{H}\!om_{\mathcal{O}_X}(\mathcal{F},\mathcal{H})\bigr)_x
  \longrightarrow
  \operatorname{Hom}_{\mathcal{O}_x}(\mathcal{F}_x,\mathcal{H}_x)
\]
is \emph{bijective} (T, 4.1.1).
\end{env}

\begin{env}[5.2.7]
\label{0.5.2.7}
Let $\mathcal{F}$, $\mathcal{G}$ be two $\mathcal{O}_X$-modules having a
finite presentation. If, for some $x\in X$, $\mathcal{F}_x$ and $\mathcal{G}_x$
are \emph{isomorphic} $\mathcal{O}_x$-modules, then there exists an open
neighbourhood $U$ of $x$ such that $\mathcal{F}|U$ and $\mathcal{G}|U$ are
\emph{isomorphic}. Indeed, if
$\varphi:\mathcal{F}_x\to\mathcal{G}_x$ and
$\psi:\mathcal{G}_x\to\mathcal{F}_x$ are two mutually inverse isomorphisms, there
exist, by (5.2.6), an open neighbourhood $V$ of $x$ and a section $u$
(resp. $v$) of
$\mathcal{H}\!om_{\mathcal{O}_X}(\mathcal{F},\mathcal{G})$ (resp.
$\mathcal{H}\!om_{\mathcal{O}_X}(\mathcal{G},\mathcal{F})$) over $V$
such

\oldpage[0\textsubscript{I}]{47}

that $u_x=\varphi$ (resp. $v_x=\psi$). Since $(u\circ v)_x$ and
$(v\circ u)_x$ are the identity automorphisms, there exists an open
neighbourhood $U\subset V$ of $x$ such that $(u\circ v)|U$ and $(v\circ u)|U$ are the
identity automorphisms, whence the proposition.
\end{env}

\subsection{Coherent sheaves}
\label{subsection:0.5.3}

\begin{env}[5.3.1]
\label{0.5.3.1}
An $\mathcal{O}_X$-module $\mathcal{F}$ is said to be \emph{coherent} if it
satisfies the following two conditions:
\begin{enumerate}
  \item[$a)$] $\mathcal{F}$ is of finite type.
  \item[$b)$] For every open set $U\subset X$, every integer $n>0$, and every
    homomorphism
    $u:\mathcal{O}_X^n|U\to\mathcal{F}|U$, the kernel of $u$ is of finite type.
\end{enumerate}
Note that these two conditions are of a \emph{local} nature.

For most of the proofs of the properties of coherent sheaves
recalled in what follows, cf. (FAC), I, 2.
\end{env}

\begin{env}[5.3.2]
\label{0.5.3.2}
Every coherent $\mathcal{O}_X$-module admits a finite presentation (5.2.5); the
converse is not necessarily true, since $\mathcal{O}_X$ itself is not
necessarily a coherent $\mathcal{O}_X$-module.

Every sub-$\mathcal{O}_X$-module \emph{of finite type} of a coherent
$\mathcal{O}_X$-module is coherent; every \emph{finite} direct sum
of coherent $\mathcal{O}_X$-modules is a coherent $\mathcal{O}_X$-module.
\end{env}

\begin{env}[5.3.3]
\label{0.5.3.3}
If $0\to\mathcal{F}\to\mathcal{G}\to\mathcal{H}\to0$ is an exact sequence of
$\mathcal{O}_X$-modules and two of these $\mathcal{O}_X$-modules are
coherent, then so is the third.
\end{env}

\begin{env}[5.3.4]
\label{0.5.3.4}
If $\mathcal{F}$ and $\mathcal{G}$ are two coherent $\mathcal{O}_X$-modules and
$u:\mathcal{F}\to\mathcal{G}$ is a homomorphism, then $\operatorname{Im}(u)$,
$\operatorname{Ker}(u)$ and $\operatorname{Coker}(u)$ are coherent
$\mathcal{O}_X$-modules. In particular, if $\mathcal{F}$ and
$\mathcal{G}$ are coherent sub-$\mathcal{O}_X$-modules of a
coherent $\mathcal{O}_X$-module, then $\mathcal{F}+\mathcal{G}$ and
$\mathcal{F}\cap\mathcal{G}$ are coherent.
\end{env}

\begin{env}[5.3.5]
\label{0.5.3.5}
If $\mathcal{F}$ and $\mathcal{G}$ are two coherent $\mathcal{O}_X$-modules,
so are $\mathcal{F}\otimes_{\mathcal{O}_X}\mathcal{G}$ and
$\mathcal{H}\!om_{\mathcal{O}_X}(\mathcal{F},\mathcal{G})$.
\end{env}

\begin{env}[5.3.6]
\label{0.5.3.6}
Let $\mathcal{F}$ be a coherent $\mathcal{O}_X$-module and $\mathcal{J}$ a
coherent sheaf of ideals of $\mathcal{O}_X$. Then the
$\mathcal{O}_X$-module $\mathcal{J}\mathcal{F}$ is coherent, as the image
of $\mathcal{J}\otimes_{\mathcal{O}_X}\mathcal{F}$ under the canonical
homomorphism
$\mathcal{J}\otimes_{\mathcal{O}_X}\mathcal{F}\to\mathcal{F}$
(5.3.4 and 5.3.5).
\end{env}

\begin{env}[5.3.7]
\label{0.5.3.7}
An $\mathcal{O}_X$-algebra $\mathcal{A}$ is said to be \emph{coherent} if
$\mathcal{A}$ is a coherent $\mathcal{O}_X$-module. In particular,
$\mathcal{O}_X$ is a \emph{coherent sheaf of rings} if and only if,
for every open set $U\subset X$ and every homomorphism of the form
$u:\mathcal{O}_X^p|U\to\mathcal{O}_X|U$, the kernel of $u$ is an
$(\mathcal{O}_X|U)$-module of finite type.

If $\mathcal{O}_X$ is a coherent sheaf of rings, every
$\mathcal{O}_X$-module $\mathcal{F}$ admitting a finite presentation (5.2.5)
is coherent, by virtue of (5.3.4).

The \emph{annihilator} of an $\mathcal{O}_X$-module $\mathcal{F}$ is the kernel
$\mathcal{J}$ of the canonical homomorphism
$\mathcal{O}_X\to\mathcal{H}\!om_{\mathcal{O}_X}(\mathcal{F},\mathcal{F})$
which assigns to every section $s\in\Gamma(U,\mathcal{O}_X)$ the
multiplication by $s$ in
$\operatorname{Hom}(\mathcal{F}|U,\mathcal{F}|U)$; if $\mathcal{O}_X$ is
coherent and $\mathcal{F}$ is a coherent $\mathcal{O}_X$-module, then
$\mathcal{J}$ is coherent (5.3.4 and 5.3.5) and, for every $x\in X$,
$\mathcal{J}_x$ is the annihilator of $\mathcal{F}_x$ (5.2.6).
\end{env}

\oldpage[0\textsubscript{I}]{48}

\begin{env}[5.3.8]
\label{0.5.3.8}
Suppose $\mathcal{O}_X$ coherent; let $\mathcal{F}$ be a coherent
$\mathcal{O}_X$-module, $x$ a point of $X$, and $M$ a submodule of
finite type of $\mathcal{F}_x$; then there exist an open neighbourhood $U$ of $x$
and a coherent sub-$(\mathcal{O}_X|U)$-module $\mathcal{G}$ of
$\mathcal{F}|U$ such that $\mathcal{G}_x=M$ (T, 4.1, lemma 1).

This result, together with the properties of the sub-$\mathcal{O}_X$-modules of a
coherent $\mathcal{O}_X$-module, imposes necessary conditions on the rings
$\mathcal{O}_x$ for $\mathcal{O}_X$ to be coherent. For example (5.3.4),
the intersection of two ideals of finite type of $\mathcal{O}_x$ must again be
an ideal of finite type.
\end{env}

\begin{env}[5.3.9]
\label{0.5.3.9}
Suppose $\mathcal{O}_X$ coherent, and let $M$ be an $\mathcal{O}_x$-module
admitting a finite presentation, hence isomorphic to a cokernel of a
homomorphism $\varphi:\mathcal{O}_x^p\to\mathcal{O}_x^q$; then there exist an
open neighbourhood $U$ of $x$ and a coherent $(\mathcal{O}_X|U)$-module
$\mathcal{F}$ such that $\mathcal{F}_x$ is isomorphic to $M$. Indeed, by virtue
of (5.2.6), there exists a section $u$ of
$\mathcal{H}\!om_{\mathcal{O}_X}(\mathcal{O}_X^p,\mathcal{O}_X^q)$ over
an open neighbourhood $U$ of $x$ such that $u_x=\varphi$; the cokernel
$\mathcal{F}$ of the homomorphism
$u:\mathcal{O}_X^p|U\to\mathcal{O}_X^q|U$ answers the question (5.3.4).
\end{env}

\begin{env}[5.3.10]
\label{0.5.3.10}
Suppose $\mathcal{O}_X$ coherent, and let $\mathcal{J}$ be a coherent sheaf
of ideals in $\mathcal{O}_X$. For an
$(\mathcal{O}_X/\mathcal{J})$-module $\mathcal{F}$ to be coherent, it is necessary and
sufficient that it be coherent as an $\mathcal{O}_X$-module. In particular,
$\mathcal{O}_X/\mathcal{J}$ is a coherent sheaf of rings.
\end{env}

\begin{env}[5.3.11]
\label{0.5.3.11}
Let $f:X\to Y$ be a morphism of ringed spaces, and suppose
$\mathcal{O}_X$ coherent; then, for every coherent $\mathcal{O}_Y$-module
$\mathcal{G}$, $f^*(\mathcal{G})$ is a coherent $\mathcal{O}_X$-module.
Indeed, with the notation of (5.2.4), one can suppose that $\mathcal{G}|V$ is
the cokernel of a homomorphism
$v:\mathcal{O}_Y^q|V\to\mathcal{O}_Y^p|V$; since $f_U^*$ is right exact,
$f^*(\mathcal{G})|U=f_U^*(\mathcal{G}|V)$ is the cokernel of the homomorphism
$f_U^*(v):\mathcal{O}_X^q|U\to\mathcal{O}_X^p|U$, whence our assertion.
\end{env}

\begin{env}[5.3.12]
\label{0.5.3.12}
Let $Y$ be a closed subset of $X$, $j:Y\to X$ the canonical injection,
$\mathcal{O}_Y$ a sheaf of rings on $Y$, and set
$\mathcal{O}_X=j_*(\mathcal{O}_Y)$. For an $\mathcal{O}_Y$-module
$\mathcal{G}$ to be of finite type (resp. quasi-coherent, coherent), it is necessary and
sufficient that $j_*(\mathcal{G})$ be an $\mathcal{O}_X$-module of finite type
(resp. quasi-coherent, coherent).
\end{env}

\subsection{Locally free sheaves}
\label{subsection:0.5.4}

\begin{env}[5.4.1]
\label{0.5.4.1}
Let $X$ be a ringed space. An $\mathcal{O}_X$-module $\mathcal{F}$
is said to be \emph{locally free} if, for every $x\in X$, there exists an open
neighbourhood $U$ of $x$ such that $\mathcal{F}|U$ is isomorphic to an
$(\mathcal{O}_X|U)$-module of the form $\mathcal{O}_X^{(I)}|U$, where $I$ may
depend on $U$. If $I$ is finite for every $U$, $\mathcal{F}$ is said to be
\emph{of finite rank}; if for every $U$, $I$ has the same finite number $n$
of elements, $\mathcal{F}$ is said to be \emph{of rank $n$}. A
locally free $\mathcal{O}_X$-module of rank 1 is also called
\emph{invertible} (cf. (5.4.3)). If $\mathcal{F}$ is a
locally free $\mathcal{O}_X$-module of finite rank, then for every $x\in X$,
$\mathcal{F}_x$ is a free $\mathcal{O}_x$-module of finite rank $n(x)$, and there
exists a neighbourhood $U$ of $x$ such that $\mathcal{F}|U$ is of rank $n(x)$;
if $X$ is connected, $n(x)$ is therefore \emph{constant}.

It is clear that every locally free sheaf is quasi-coherent, and if
$\mathcal{O}_X$ is a coherent sheaf of rings, every locally free
$\mathcal{O}_X$-module of finite rank is coherent.

If $\mathcal{L}$ is locally free,
$\mathcal{L}\otimes_{\mathcal{O}_X}\mathcal{F}$ is an \emph{exact} functor
in $\mathcal{F}$ on the category of $\mathcal{O}_X$-modules.

We shall mostly have to consider locally free $\mathcal{O}_X$-modules
of finite rank,

\oldpage[0\textsubscript{I}]{49}

and when we speak of locally free sheaves without further specification, it
will be understood that they are of \emph{finite rank}.
\end{env}

\begin{env}[5.4.2]
\label{0.5.4.2}
If $\mathcal{L}$, $\mathcal{F}$ are two $\mathcal{O}_X$-modules, there is a
canonical functorial homomorphism
\[
  \check{\mathcal{L}}\otimes_{\mathcal{O}_X}\mathcal{F}
  =\mathcal{H}\!om_{\mathcal{O}_X}(\mathcal{L},\mathcal{O}_X)
   \otimes_{\mathcal{O}_X}\mathcal{F}
  \longrightarrow
  \mathcal{H}\!om_{\mathcal{O}_X}(\mathcal{L},\mathcal{F})
  \tag{5.4.2.1}\label{0.5.4.2.1}
\]
defined as follows: for every open set $U$, to every pair $(u,t)$, where
\[
  u\in\Gamma\bigl(U,
    \mathcal{H}\!om_{\mathcal{O}_X}(\mathcal{L},\mathcal{O}_X)\bigr)
  =\operatorname{Hom}(\mathcal{L}|U,\mathcal{O}_X|U)
  \quad\text{and}\quad
  t\in\Gamma(U,\mathcal{F}),
\]
one assigns the element of
$\operatorname{Hom}(\mathcal{L}|U,\mathcal{F}|U)$ which, for every $x\in U$,
assigns to $s_x\in\mathcal{L}_x$ the element $u_x(s_x)t_x$ of
$\mathcal{F}_x$. If $\mathcal{L}$ is \emph{locally free of finite rank},
this homomorphism is \emph{bijective}; indeed, the property being local, one may
restrict oneself to the case where $\mathcal{L}=\mathcal{O}_X^n$; since, for every
$\mathcal{O}_X$-module $\mathcal{G}$,
$\mathcal{H}\!om_{\mathcal{O}_X}(\mathcal{O}_X^n,\mathcal{G})$ is
canonically isomorphic to $\mathcal{G}^n$, one is reduced to the case
$\mathcal{L}=\mathcal{O}_X$, which is immediate.
\end{env}

\begin{env}[5.4.3]\label{0.5.4.3}
\EGAReaderTarget{canon-ref.0.5.4.3-fr}{0.5.4.3-fr}
If $\mathcal{L}$ is invertible, so is its dual
$\check{\mathcal{L}}=
\mathcal{H}\!om_{\mathcal{O}_X}(\mathcal{L},\mathcal{O}_X)$, since one is
immediately reduced (the question being local) to the case $\mathcal{L}=\mathcal{O}_X$.
Moreover, there is a canonical isomorphism
\[
  \mathcal{H}\!om_{\mathcal{O}_X}(\mathcal{L},\mathcal{O}_X)
  \otimes_{\mathcal{O}_X}\mathcal{L}
  \xrightarrow{\sim}
  \mathcal{O}_X
  \tag{5.4.3.1}\label{0.5.4.3.1}
\]
indeed, by (5.4.2), it suffices to define a canonical isomorphism
$\mathcal{H}\!om_{\mathcal{O}_X}(\mathcal{L},\mathcal{L})
\xrightarrow{\sim}\mathcal{O}_X$. Now, for \emph{every}
$\mathcal{O}_X$-module $\mathcal{F}$, there is a canonical homomorphism
$\mathcal{O}_X\longrightarrow
\mathcal{H}\!om_{\mathcal{O}_X}(\mathcal{F},\mathcal{F})$ (5.3.7). It remains
to prove that if $\mathcal{F}=\mathcal{L}$ is invertible, this homomorphism
is bijective, and since the question is local, one is reduced to the case
$\mathcal{L}=\mathcal{O}_X$, which is immediate.

In view of the foregoing, we set
$\mathcal{L}^{-1}=
\mathcal{H}\!om_{\mathcal{O}_X}(\mathcal{L},\mathcal{O}_X)$ and say that
$\mathcal{L}^{-1}$ is the \emph{inverse} of $\mathcal{L}$. The terminology
``invertible sheaf'' can be justified as follows when $X$ is
reduced to a point and $\mathcal{O}_X$ is a \emph{local} ring $A$ with
maximal ideal $\mathfrak m$; if $M$ and $M'$ are two $A$-modules ($M$ being of
finite type) such that $M\otimes_A M'$ is isomorphic to $A$, then, since
$(A/\mathfrak m)\otimes_A(M\otimes_A M')$ is identified with
$(M/\mathfrak mM)\otimes_{A/\mathfrak m}(M'/\mathfrak mM')$, this last
tensor product of vector spaces over the field $A/\mathfrak m$ is
isomorphic to $A/\mathfrak m$, which requires $M/\mathfrak mM$ and
$M'/\mathfrak mM'$ to be of dimension 1. For every element $z\in M$
not belonging to $\mathfrak mM$, one therefore has $M=Az+\mathfrak mM$, which
implies $M=Az$ by Nakayama's lemma, $M$ being of finite type.
Moreover, since the annihilator of $z$ annihilates $M\otimes_A M'$, which is isomorphic to
$A$, this annihilator is $\{0\}$, and $M$ is consequently \emph{isomorphic to $A$}.
In the general case, this shows that if $\mathcal{L}$ is an
$\mathcal{O}_X$-module of finite type such that there exists an
$\mathcal{O}_X$-module $\mathcal{F}$ for which
$\mathcal{L}\otimes_{\mathcal{O}_X}\mathcal{F}$ is isomorphic to
$\mathcal{O}_X$, and if moreover the rings $\mathcal{O}_x$ are local
rings, then $\mathcal{L}_x$ is an $\mathcal{O}_x$-module isomorphic to
$\mathcal{O}_x$ for every $x\in X$. If $\mathcal{O}_X$ and $\mathcal{L}$ are
assumed to be \emph{coherent}, one therefore concludes that $\mathcal{L}$ is invertible,
by virtue of (5.2.7).
\end{env}

\begin{env}[5.4.4]
\label{0.5.4.4}
If $\mathcal{L}$ and $\mathcal{L}'$ are two invertible $\mathcal{O}_X$-modules,
so is $\mathcal{L}\otimes_{\mathcal{O}_X}\mathcal{L}'$, since, the question being
local, one may suppose that $\mathcal{L}=\mathcal{O}_X$, and the result is
then trivial. For every integer $n\geq1$, one denotes by
$\mathcal{L}^{\otimes n}$ the tensor product of $n$ sheaves

\oldpage[0\textsubscript{I}]{50}

identical to $\mathcal{L}$; by convention one sets
$\mathcal{L}^{\otimes0}=\mathcal{O}_X$ and, for $n\geq1$,
$\mathcal{L}^{\otimes(-n)}=(\mathcal{L}^{-1})^{\otimes n}$. With this
notation, there then exists a \emph{canonical functorial isomorphism}
\[
  \mathcal{L}^{\otimes m}\otimes_{\mathcal{O}_X}\mathcal{L}^{\otimes n}
  \xrightarrow{\sim}
  \mathcal{L}^{\otimes(n+m)}
  \tag{5.4.4.1}\label{0.5.4.4.1}
\]
for all rational integers $m,n$: indeed, by virtue of the
definitions, one is immediately reduced to the case where $m=-1$, $n=1$, and
the isomorphism in question was then defined in (5.4.3).
\end{env}

\begin{env}[5.4.5]\label{0.5.4.5}
\EGAReaderTarget{canon-ref.0.5.4.5-fr}{0.5.4.5-fr}
Let $f:Y\to X$ be a morphism of ringed spaces. If $\mathcal{L}$ is a
locally free (resp. invertible) $\mathcal{O}_X$-module, then
$f^*(\mathcal{L})$ is a locally free (resp.
invertible) $\mathcal{O}_Y$-module: this follows at once from the facts that the inverse images of two
locally isomorphic $\mathcal{O}_X$-modules are locally isomorphic, that
$f^*$ commutes with finite direct sums, and that
$f^*(\mathcal{O}_X)=\mathcal{O}_Y$ (4.3.4). Moreover, we know that there is a
canonical functorial homomorphism
\[
  f^*(\check{\mathcal{L}})
  \longrightarrow
  \check{\bigl(f^*(\mathcal{L})\bigr)}
\]
(4.4.6), and when $\mathcal{L}$ is locally free, this homomorphism is
\emph{bijective}: indeed, one is again reduced to the case where
$\mathcal{L}=\mathcal{O}_X$, which is trivial. One concludes that if
$\mathcal{L}$ is invertible, $f^*(\mathcal{L}^{\otimes n})$ is canonically
identified with $(f^*(\mathcal{L}))^{\otimes n}$ for every rational integer
$n$.
\end{env}

\begin{env}[5.4.6]
\label{0.5.4.6}
Let $\mathcal{L}$ be an invertible $\mathcal{O}_X$-module; one denotes by
$\Gamma_*(X,\mathcal{L})$, or simply $\Gamma_*(\mathcal{L})$, the abelian
group that is the direct sum
$\bigoplus_{n\in\mathbb{Z}}\Gamma(X,\mathcal{L}^{\otimes n})$; it is endowed
with a structure of \emph{graded ring} by assigning to the pair
$(s_n,s_m)$, where $s_n\in\Gamma(X,\mathcal{L}^{\otimes n})$,
$s_m\in\Gamma(X,\mathcal{L}^{\otimes m})$, the section of
$\mathcal{L}^{\otimes(n+m)}$ over $X$ that corresponds canonically
(5.4.4.1) to the section $s_n\otimes s_m$ of
$\mathcal{L}^{\otimes n}\otimes_{\mathcal{O}_X}\mathcal{L}^{\otimes m}$;
the associativity of this multiplication is verified immediately. It is
clear that $\Gamma_*(X,\mathcal{L})$ is a covariant functor in
$\mathcal{L}$, with values in the category of graded rings.

If now $\mathcal{F}$ is an arbitrary $\mathcal{O}_X$-module, one sets
\[
  \Gamma_*(\mathcal{L},\mathcal{F})
  =\bigoplus_{n\in\mathbb{Z}}
   \Gamma\bigl(X,
     \mathcal{F}\otimes_{\mathcal{O}_X}\mathcal{L}^{\otimes n}\bigr).
\]
This abelian group is endowed with a structure of \emph{graded module} over
the graded ring $\Gamma_*(\mathcal{L})$ as follows: to the pair
$(s_n,u_m)$, where $s_n\in\Gamma(X,\mathcal{L}^{\otimes n})$ and
$u_m\in\Gamma(X,
\mathcal{F}\otimes_{\mathcal{O}_X}\mathcal{L}^{\otimes m})$, one
assigns the section of
$\mathcal{F}\otimes_{\mathcal{O}_X}\mathcal{L}^{\otimes(m+n)}$ that corresponds
canonically (5.4.4.1) to $s_n\otimes u_m$; the verification of the module
axioms is immediate. For fixed $X$ and $\mathcal{L}$,
$\Gamma_*(\mathcal{L},\mathcal{F})$ is a covariant functor in
$\mathcal{F}$ with values in the category of graded
$\Gamma_*(\mathcal{L})$-modules; for fixed $X$ and $\mathcal{F}$,
it is a covariant functor in $\mathcal{L}$ with values in the category of
abelian groups.

If $f:Y\to X$ is a morphism of ringed spaces, the canonical homomorphism
(4.4.3.2)
\[
  \rho:\mathcal{L}^{\otimes n}
  \longrightarrow
  f_*\bigl(f^*(\mathcal{L}^{\otimes n})\bigr)
\]
defines a homomorphism of abelian groups
\[
  \Gamma(X,\mathcal{L}^{\otimes n})
  \longrightarrow
  \Gamma\bigl(Y,f^*(\mathcal{L}^{\otimes n})\bigr),
\]
and since
$f^*(\mathcal{L}^{\otimes n})=(f^*(\mathcal{L}))^{\otimes n}$, it follows from the
definitions of the canonical homomorphisms (4.4.3.2) and (5.4.4.1) that the
preceding homomorphisms define a \emph{functorial homomorphism
of graded rings}
$\Gamma_*(\mathcal{L})\to\Gamma_*(f^*(\mathcal{L}))$. The same canonical
homomorphism (4.4.3) likewise defines a homomorphism of abelian groups
\[
  \Gamma\bigl(X,
    \mathcal{F}\otimes_{\mathcal{O}_X}\mathcal{L}^{\otimes n}\bigr)
  \longrightarrow
  \Gamma\bigl(Y,
    f^*(\mathcal{F}\otimes_{\mathcal{O}_X}\mathcal{L}^{\otimes n})\bigr),
\]
and since
\[
  f^*(\mathcal{F}\otimes_{\mathcal{O}_X}\mathcal{L}^{\otimes n})
  =f^*(\mathcal{F})\otimes_{\mathcal{O}_Y}
   (f^*(\mathcal{L}))^{\otimes n}
  \tag{4.3.3.1}
\]

\oldpage[0\textsubscript{I}]{51}

these homomorphisms (for variable $n$) define a
\emph{di-homomorphism of graded modules}
\[
  \Gamma_*(\mathcal{L},\mathcal{F})
  \longrightarrow
  \Gamma_*(f^*(\mathcal{L}),f^*(\mathcal{F})).
\]
\end{env}

\begin{env}[5.4.7]\label{0.5.4.7}
\EGAReaderTarget{canon-ref.0.5.4.7-fr}{0.5.4.7-fr}
One can show that there exists a \emph{set} $\mathfrak{M}$ (also denoted
$\mathfrak{M}(X)$) of invertible $\mathcal{O}_X$-modules such that every
invertible $\mathcal{O}_X$-module is isomorphic to one and only one element of
$\mathfrak{M}$\footnote{See the book in preparation cited in
the Introduction.}; one defines on $\mathfrak{M}$ a law of composition by
assigning to two elements $\mathcal{L}$, $\mathcal{L}'$ of
$\mathfrak{M}$ the unique element of $\mathfrak{M}$ isomorphic to
$\mathcal{L}\otimes_{\mathcal{O}_X}\mathcal{L}'$. With this law of
composition, $\mathfrak{M}$ is a \emph{group isomorphic to the cohomology
group} $H^1(X,\mathcal{O}_X^*)$, where $\mathcal{O}_X^*$ is the
sub-sheaf of $\mathcal{O}_X$ such that $\Gamma(U,\mathcal{O}_X^*)$ is the
group of invertible elements of the ring $\Gamma(U,\mathcal{O}_X)$ for
every open set $U\subset X$ (so $\mathcal{O}_X^*$ is a sheaf of
\emph{multiplicative} abelian groups).

To see this, note that for every open set $U\subset X$, the group of sections
$\Gamma(U,\mathcal{O}_X^*)$ is canonically identified with the \emph{group of
automorphisms} of the $(\mathcal{O}_X|U)$-module $\mathcal{O}_X|U$,
the identification assigning to a section $\varepsilon$ of
$\mathcal{O}_X^*$ over $U$ the automorphism $u$ of $\mathcal{O}_X|U$
such that $u_x(s_x)=\varepsilon_xs_x$ for every $x\in U$ and every
$s_x\in\mathcal{O}_x$. Now let $\mathfrak{U}=(U_\lambda)$ be an open
covering of $X$; giving, for every pair of indices $(\lambda,\mu)$, an
automorphism $\theta_{\lambda\mu}$ of
$\mathcal{O}_X|(U_\lambda\cap U_\mu)$ amounts to giving a
$1$-\emph{cochain} of the covering $\mathfrak{U}$ with values in
$\mathcal{O}_X^*$, and to say that the $\theta_{\lambda\mu}$ satisfy the
gluing condition (3.3.1) means that the corresponding cochain is
a \emph{cocycle}. Likewise, giving, for every $\lambda$, an automorphism
$\omega_\lambda$ of $\mathcal{O}_X|U_\lambda$ amounts to giving a
$0$-cochain of the covering $\mathfrak{U}$ with values in
$\mathcal{O}_X^*$, and its \emph{coboundary} corresponds to the family of
automorphisms
$(\omega_\lambda|U_\lambda\cap U_\mu)\circ
 (\omega_\mu|U_\lambda\cap U_\mu)^{-1}$.

To every $1$-cocycle of $\mathfrak{U}$ with values
in $\mathcal{O}_X^*$ one can assign the element of $\mathfrak{M}$ isomorphic to the
invertible $\mathcal{O}_X$-module obtained by gluing by means of the
family of automorphisms $(\theta_{\lambda\mu})$ corresponding to this cocycle,
and to two cohomologous cocycles there correspond two equal elements of
$\mathfrak{M}$ (3.3.2); in other words, one has thus defined a map
$\varphi_{\mathfrak{U}}:H^1(\mathfrak{U},\mathcal{O}_X^*)\to\mathfrak{M}$.
Moreover, if $\mathfrak{B}$ is a second open covering of $X$, finer
than $\mathfrak{U}$, the diagram
\[
\begin{tikzcd}[column sep=huge,row sep=large]
H^1(\mathfrak{U},\mathcal{O}_X^*)
  \arrow[dd]
  \arrow[rrd,"\varphi_{\mathfrak{U}}"] & & \\
& & \mathfrak{M} \\
H^1(\mathfrak{B},\mathcal{O}_X^*)
  \arrow[rru,swap,"\varphi_{\mathfrak{B}}"] & &
\end{tikzcd}
\]
where the vertical arrow is the canonical homomorphism (G, II, 5.7), is
commutative, as follows from (3.3.3). By passage to the inductive limit, one
therefore indeed obtains a map
$H^1(X,\mathcal{O}_X^*)\to\mathfrak{M}$, the \v{C}ech cohomology group
$\check{H}^1(X,\mathcal{O}_X^*)$ being identified, as is well known, with the first
cohomology group $H^1(X,\mathcal{O}_X^*)$ (G, II, 5.9, cor. of th.~5.9.1).
This map is \emph{surjective}: indeed, by definition, for every invertible
$\mathcal{O}_X$-module $\mathcal{L}$, there is an open covering
$(U_\lambda)$ of $X$ such that $\mathcal{L}$ is obtained by gluing the
sheaves $\mathcal{O}_X|U_\lambda$ (3.3.1). It is also
\emph{injective}, for it suffices to prove this for the maps
$H^1(\mathfrak{U},\mathcal{O}_X^*)\to\mathfrak{M}$, and this then follows from
(3.3.2). It remains to show that

\oldpage[0\textsubscript{I}]{52}

the bijection thus defined is a group homomorphism. Now, given
two invertible $\mathcal{O}_X$-modules $\mathcal{L}$, $\mathcal{L}'$, there is
an open covering $(U_\lambda)$ such that $\mathcal{L}|U_\lambda$ and
$\mathcal{L}'|U_\lambda$ are isomorphic to $\mathcal{O}_X|U_\lambda$ for
every $\lambda$; there is therefore, for each index $\lambda$, an element
$a_\lambda$ (resp. $a'_\lambda$) of $\Gamma(U_\lambda,\mathcal{L})$ (resp.
$\Gamma(U_\lambda,\mathcal{L}'))$ such that the elements of
$\Gamma(U_\lambda,\mathcal{L})$ (resp. $\Gamma(U_\lambda,\mathcal{L}'))$
are the $s_\lambda\cdot a_\lambda$ (resp.
$s_\lambda\cdot a'_\lambda$), where $s_\lambda$ ranges over
$\Gamma(U_\lambda,\mathcal{O}_X)$. The corresponding cocycles
$(\varepsilon_{\lambda\mu})$, $(\varepsilon'_{\lambda\mu})$ are such that
$s_\lambda\cdot a_\lambda=s_\mu\cdot a_\mu$ (resp.
$s_\lambda\cdot a'_\lambda=s_\mu\cdot a'_\mu$) over
$U_\lambda\cap U_\mu$ is equivalent to
$s_\lambda=\varepsilon_{\lambda\mu}s_\mu$ (resp.
$s_\lambda=\varepsilon'_{\lambda\mu}s_\mu$) over
$U_\lambda\cap U_\mu$. Since the sections of
$\mathcal{L}\otimes_{\mathcal{O}_X}\mathcal{L}'$ over $U_\lambda$
are the finite sums of the
$s_\lambda s'_\lambda\cdot(a_\lambda\otimes a'_\lambda)$, where $s_\lambda$ and
$s'_\lambda$ range over $\Gamma(U_\lambda,\mathcal{O}_X)$, it is clear that
the cocycle $(\varepsilon_{\lambda\mu}\varepsilon'_{\lambda\mu})$ corresponds
to $\mathcal{L}\otimes_{\mathcal{O}_X}\mathcal{L}'$, which completes the
proof\footnote{For a general form of this result, see the book
cited in the footnote on p.~51.}.
\end{env}

\begin{env}[5.4.8]
\label{0.5.4.8}
Let $f=(\psi,\omega)$ be a morphism $Y\to X$ of ringed spaces. The functor
$f^*$, restricted to the category of invertible $\mathcal{O}_X$-modules,
defines a map (again denoted by $f^*$, by abuse of language) from the set
$\mathfrak{M}(X)$ to the set $\mathfrak{M}(Y)$. On the other hand, there is a
canonical homomorphism (T, 3.2.2)
\[
  H^1(X,\mathcal{O}_X^*)\longrightarrow H^1(Y,\mathcal{O}_Y^*).
  \tag{5.4.8.1}\label{0.5.4.8.1}
\]
When $\mathfrak{M}(X)$ and
$H^1(X,\mathcal{O}_X^*)$ (resp. $\mathfrak{M}(Y)$ and
$H^1(Y,\mathcal{O}_Y^*)$) are canonically identified (5.4.7), the homomorphism (5.4.8.1) \emph{is identified with
the map $f^*$}. Indeed, if $\mathcal{L}$ comes from a cocycle
$(\varepsilon_{\lambda\mu})$ corresponding to an open covering
$(U_\lambda)$ of $X$, it suffices to show that $f^*(\mathcal{L})$ comes
from a cocycle whose cohomology class is the image under (5.4.8.1) of that of
$(\varepsilon_{\lambda\mu})$. Now, if $\theta_{\lambda\mu}$ is
the automorphism of $\mathcal{O}_X|(U_\lambda\cap U_\mu)$ corresponding to
$\varepsilon_{\lambda\mu}$, it is clear that $f^*(\mathcal{L})$ is obtained by
gluing the $\mathcal{O}_Y|\psi^{-1}(U_\lambda)$ by means of the
automorphisms $f^*(\theta_{\lambda\mu})$, and it therefore suffices to verify that
the latter correspond to the cocycle
$(\omega^\#(\varepsilon_{\lambda\mu}))$, which follows at once from the
definitions (we have identified $\varepsilon_{\lambda\mu}$ with its canonical image
under $\rho$ (3.7.2), a section of $\psi^*(\mathcal{O}_X^*)$ over
$\psi^{-1}(U_\lambda\cap U_\mu)$).
\end{env}

\begin{env}[5.4.9]
\label{0.5.4.9}
Let $\mathcal{E}$, $\mathcal{F}$ be two $\mathcal{O}_X$-modules,
$\mathcal{F}$ being assumed \emph{locally free}, and let $\mathcal{G}$ be an
$\mathcal{O}_X$-module that is an \emph{extension of $\mathcal{F}$ by $\mathcal{E}$},
in other words such that there exists an exact sequence
\[
  0\longrightarrow\mathcal{E}\xrightarrow{i}\mathcal{G}
  \xrightarrow{p}\mathcal{F}\longrightarrow 0.
\]
Then, for every $x\in X$, there exists an open neighbourhood $U$ of $x$ such that
$\mathcal{G}|U$ is isomorphic to the \emph{direct sum}
$\mathcal{E}|U\oplus\mathcal{F}|U$. Indeed, one may restrict oneself to the case where
$\mathcal{F}=\mathcal{O}_X^n$; let $e_i$ ($1\leq i\leq n$) be the canonical
sections (5.5.5) of $\mathcal{O}_X^n$; there then exist an open neighbourhood
$U$ of $x$ and $n$ sections $s_i$ of $\mathcal{G}$ over $U$ such that
$p(s_i|U)=e_i|U$ for $1\leq i\leq n$. That being so, let $f$ be
the homomorphism $\mathcal{F}|U\to\mathcal{G}|U$ defined by the sections
$s_i|U$ (5.1.1). It is immediate that for every open set $V\subset U$ and every
section $s\in\Gamma(V,\mathcal{G})$ one has
$s-f(p(s))\in\Gamma(V,\mathcal{E})$, whence our assertion.
\end{env}

\begin{env}[5.4.10]
\label{0.5.4.10}
Let $f:X\to Y$ be a morphism of ringed spaces, $\mathcal{F}$ an
$\mathcal{O}_X$-module, and $\mathcal{L}$ a locally free
$\mathcal{O}_Y$-module of finite rank. There then exists a canonical isomorphism
\[
  f_*(\mathcal{F})\otimes_{\mathcal{O}_Y}\mathcal{L}
  \xrightarrow{\sim}
  f_*\bigl(\mathcal{F}\otimes_{\mathcal{O}_X}f^*(\mathcal{L})\bigr).
  \tag{5.4.10.1}\label{0.5.4.10.1}
\]

\oldpage[0\textsubscript{I}]{53}

Indeed, for every $\mathcal{O}_Y$-module $\mathcal{L}$, there is a
canonical homomorphism
\[
  f_*(\mathcal{F})\otimes_{\mathcal{O}_Y}\mathcal{L}
  \xrightarrow{1\otimes\rho}
  f_*(\mathcal{F})\otimes_{\mathcal{O}_Y}f_*(f^*(\mathcal{L}))
  \xrightarrow{\alpha}
  f_*\bigl(\mathcal{F}\otimes_{\mathcal{O}_X}f^*(\mathcal{L})\bigr),
\]
where $\rho$ is the homomorphism (4.4.3.2) and $\alpha$ the homomorphism (4.2.2.1).
To show that this homomorphism is bijective when $\mathcal{L}$ is locally free,
it suffices, the question being local on $Y$, to consider the
case where $\mathcal{L}=\mathcal{O}_Y^n$; moreover, since $f_*$ and $f^*$ commute with
finite direct sums, one may suppose $n=1$, and in this case the proposition
follows at once from the definitions and from the relation
$f^*(\mathcal{O}_Y)=\mathcal{O}_X$.
\end{env}


\subsection{Sheaves on a ringed space in local rings}
\label{subsection:0.5.5}

\begin{env}[5.5.1]
\label{0.5.5.1}
We shall say that a ringed space $(X,\mathcal{O}_X)$ is a \emph{ringed
space in local rings} if, for every $x\in X$, $\mathcal{O}_x$ is a
local ring; these ringed spaces will be by far the most frequent of those
considered in this work. We shall then denote by
$\mathfrak{m}_x$ the \emph{maximal ideal} of $\mathcal{O}_x$, and by $k(x)$ the
\emph{residue field} $\mathcal{O}_x/\mathfrak{m}_x$; for every
$\mathcal{O}_X$-module $\mathcal{F}$, every open set $U$ of $X$, every point
$x\in U$ and every section $f\in\Gamma(U,\mathcal{F})$, we shall denote by
$f(x)$ the \emph{class} of the germ
$f_x\in\mathcal{F}_x\pmod{\mathfrak{m}_x\mathcal{F}_x}$, and we shall say that
it is the \emph{value} of $f$ at the point $x$. The relation $f(x)=0$ thus
means that $f_x\in\mathfrak{m}_x\mathcal{F}_x$; when it holds, we shall say
(by abuse of language) that $f$ \emph{vanishes at $x$}. Care must be taken not to
confuse this relation with $f_x=0$.
\end{env}

\begin{env}[5.5.2]
\label{0.5.5.2}
Let $X$ be a ringed space in local rings, $\mathcal{L}$ an
invertible $\mathcal{O}_X$-module, and $f$ a section of $\mathcal{L}$ over
$X$. Then the following three properties at a point $x\in X$ are
\emph{equivalent}:
\begin{enumerate}
  \item[$a)$] \emph{$f_x$ is a generator of $\mathcal{L}_x$};
  \item[$b)$] \emph{$f_x\notin\mathfrak{m}_x\mathcal{L}_x$} (in other words,
    $f(x)\neq0$);
  \item[$c)$] \emph{There exists a section $g$ of $\mathcal{L}^{-1}$ over
    an open neighbourhood $V$ of $x$ such that the canonical image of $f\otimes
    g$ in $\Gamma(V,\mathcal{O}_X)$ (5.4.3) is the unit section}.
\end{enumerate}

Indeed, the question being local, one may restrict oneself to the case where
$\mathcal{L}=\mathcal{O}_X$; the equivalence of $a)$ and $b)$ is then
obvious, and it is clear that $c)$ implies $b)$. On the other hand, if
$f_x\notin\mathfrak{m}_x$, then $f_x$ is invertible in $\mathcal{O}_x$, say
$f_xg_x=1_x$. By the definition of germs of sections, this means that there
exist a neighbourhood $V$ of $x$ and a section $g$ of $\mathcal{O}_X$ over
$V$ such that $fg=1$ in $V$, whence $c)$.

It follows at once from condition $c)$ that the set $X_f$ of the $x$
satisfying the equivalent conditions $a)$, $b)$, $c)$ is \emph{open} in
$X$; following the terminology introduced in (5.5.1), it is the set of the
$x$ at which $f$ \emph{does not vanish}.
\end{env}

\begin{env}[5.5.3]
\label{0.5.5.3}
Under the hypotheses of (5.5.2), let $\mathcal{L}'$ be a second invertible
$\mathcal{O}_X$-module; then, if
$f\in\Gamma(X,\mathcal{L})$ and $g\in\Gamma(X,\mathcal{L}')$, one has
\[
  X_f\cap X_g=X_{f\otimes g}.
\]
Indeed, one is immediately reduced to the case where
$\mathcal{L}=\mathcal{L}'=\mathcal{O}_X$ (the question being local); since
$f\otimes g$ is then canonically identified with the product $fg$, the proposition
is obvious.
\end{env}

\oldpage[0\textsubscript{I}]{54}

\begin{env}[5.5.4]
\label{0.5.5.4}
Let $\mathcal{F}$ be a locally free $\mathcal{O}_X$-module of rank $n$;
it is immediate that $\bigwedge^p\mathcal{F}$ is a locally free
$\mathcal{O}_X$-module of rank $\binom{n}{p}$ if $p\leq n$, and reduced to $0$ if $p>n$,
since the question is local and one is reduced to the case where
$\mathcal{F}=\mathcal{O}_X^n$; moreover, for every $x\in X$,
\[
  (\bigwedge^p\mathcal{F})_x/
  \mathfrak{m}_x(\bigwedge^p\mathcal{F})_x
\]
is a vector space of dimension $\binom{n}{p}$ over $k(x)$, which
is canonically identified with
\[
  \bigwedge^p(\mathcal{F}_x/\mathfrak{m}_x\mathcal{F}_x).
\]
Let
$s_1,\ldots,s_p$ be sections of $\mathcal{F}$ over an open set $U$ of $X$, and
let $s=s_1\wedge\cdots\wedge s_p$, which is a section of
$\bigwedge^p\mathcal{F}$ over $U$ (4.1.5); one has
$s(x)=s_1(x)\wedge\cdots\wedge s_p(x)$, and consequently, to say that
$s_1(x),\ldots,s_p(x)$ are \emph{linearly dependent} means that
$s(x)=0$. One concludes that \emph{the set of $x\in U$ such that
$s_1(x),\ldots,s_p(x)$ are linearly independent is open in
$X$}: indeed, it suffices, reducing to the case where
$\mathcal{F}=\mathcal{O}_X^n$, to apply (5.5.2) to the section that is the image of $s$
under one of the projections of
$\bigwedge^p\mathcal{F}=\mathcal{O}_X^{\binom{n}{p}}$ onto the
$\binom{n}{p}$ factors.

In particular, if $s_1,\ldots,s_n$ are $n$ sections of $\mathcal{F}$
over $U$ such that $s_1(x),\ldots,s_n(x)$ are linearly
independent at every point $x\in U$, the homomorphism
$u:\mathcal{O}_X^n|U\to\mathcal{F}|U$ defined by the $s_i$ (5.1.1) is an
\emph{isomorphism}: indeed, one may restrict oneself to the case where
$\mathcal{F}=\mathcal{O}_X^n$ and where
$\bigwedge^n\mathcal{F}$ and $\mathcal{O}_X$ are canonically identified;
$s=s_1\wedge\cdots\wedge s_n$ is then an \emph{invertible} section of $\mathcal{O}_X$
over $U$, and one defines an inverse homomorphism
of $u$ by means of Cramer's formulas.
\end{env}

\begin{env}[5.5.5]
\label{0.5.5.5}
Let $\mathcal{E}$, $\mathcal{F}$ be two locally free $\mathcal{O}_X$-modules
(of finite rank), and let $u:\mathcal{E}\to\mathcal{F}$ be a
homomorphism. For there to exist a neighbourhood $U$ of $x\in X$ such that $u|U$
is \emph{injective} and $\mathcal{F}|U$ is \emph{the direct sum of
$u(\mathcal{E})|U$ and a locally free sub-$(\mathcal{O}_X|U)$-module
$\mathcal{G}$}, it is necessary and sufficient that
$u_x:\mathcal{E}_x\to\mathcal{F}_x$ give, by passage to the quotients, an
\emph{injective} homomorphism of vector spaces
\[
  \mathcal{E}_x/\mathfrak{m}_x\mathcal{E}_x
  \longrightarrow
  \mathcal{F}_x/\mathfrak{m}_x\mathcal{F}_x.
\]
The condition is indeed \emph{necessary}, since $\mathcal{F}_x$ is then
the direct sum of the free $\mathcal{O}_x$-modules $u_x(\mathcal{E}_x)$ and
$\mathcal{G}_x$, so that $\mathcal{F}_x/\mathfrak{m}_x\mathcal{F}_x$ is the
direct sum of
$u_x(\mathcal{E}_x)/\mathfrak{m}_x u_x(\mathcal{E}_x)$ and
$\mathcal{G}_x/\mathfrak{m}_x\mathcal{G}_x$. The condition is
\emph{sufficient}, for one may restrict oneself to the case where
$\mathcal{E}=\mathcal{O}_X^m$; let $s_1,\ldots,s_m$ be the images under $u$ of the
sections $e_i$ of $\mathcal{O}_X^m$ such that $(e_i)_y$ is equal to the $i$-th
element of the canonical basis of $\mathcal{O}_y^m$ for every $y\in X$
(\emph{canonical sections of $\mathcal{O}_X^m$}); by hypothesis
$s_1(x),\ldots,s_m(x)$ are linearly independent, so, if $\mathcal{F}$
is of rank $n$, there exist $n-m$ sections $s_{m+1},\ldots,s_n$ of
$\mathcal{F}$ over a neighbourhood $V$ of $x$ such that the $s_i(x)$
$(1\leq i\leq n)$ form a basis of
$\mathcal{F}_x/\mathfrak{m}_x\mathcal{F}_x$. There then exists (5.5.4) a
neighbourhood $U\subset V$ of $x$ such that the $s_i(y)$ $(1\leq i\leq n)$ form
a basis of $\mathcal{F}_y/\mathfrak{m}_y\mathcal{F}_y$ for every $y\in U$,
and one concludes (5.5.4) that there is an isomorphism of $\mathcal{F}|U$ onto
$\mathcal{O}_X^n|U$ mapping the $s_i|U$ $(1\leq i\leq m)$ to the
$e_i|U$, which completes the proof.
\end{env}

